
\AddSymbolsB
\pagecolor{cyan!10!gray}
\color{white}
\quad\quad \lettrine[,lraise=0.2,lhang=0.5]{\textbf{伊}}{藤}佐一彻底地安下心来；他也有模有样地学着教授之前的状态，假装耐心地等待着那两杯拿铁的到来。

咖啡师们马不停蹄地在咖啡机之间奔走着——没有人比他们更忙了。伊藤佐一强装镇定，可是内心的不安让他的手指开始骚动起来。

他回头看看教授；教授似乎比往常虚弱得多，蜷缩着身子，背后披着一件红色格子衫。伊藤佐一惆怅地看着咖啡台，看着那不锈钢表面映出的模糊的倒影。

\twkkai{两杯拿铁，请慢用！}

伊藤佐一的手早已在等待过程中冻僵了；他连忙抓起两杯拿铁，走回座位。

\twkkai{教授，这是你的拿铁……你的状态似乎不太好呀……；}伊藤佐一探出头问着。

\twkkai{是啊，是啊；我也不知道怎么回事……主要是很烦……家里的东西全被弄乱了……弄得我现在神经兮兮的……}

\twkkai{嗯……是进贼了吗？}

\twkkai{不不……是我自己捡的流浪猫闯祸了——后来我把它赶出去了；}教授摇了摇头，又用手抱着头。

\twkkai{猫？}伊藤佐一的担忧加剧了。

……

课室内热闹腾腾的。只有铃木和子安静地趴在座位上。她拿着手里的奖章，感到惴惴不安。奖章的多次出现，早已让她变得极为警惕——夜晚，她总是辗转反侧，看着明媚的月光，心中泛起无限波澜。

铃声响起，课室安静下来了；坐在铃木和子隔壁的那位社团同学推了推她。她终于抬起了头，而后观察着。伊藤佐一依旧没有来；可是教授此时也并没有在场。她看了看手表。

\twkkai{今天不是星期一吗？}她思忖着。她再次看向奖章的铭文——那个她始终没有头绪、需要上下文才能理解的问题：\[\text{\twkkai{证明}}L(1,\chi)\neq 0\quad(\chi\neq 1)\]\index{ly@$L(1,\chi)\neq 0$（非平凡特征非零性）}
\quad\quad 就这样，过了约莫三分钟后，课室内再次响起欢乐的笑声。可是没一会儿，传来一阵严厉的呵斥。

\twkkai{全体安静！}

一只猫走了进来，跳上了讲台，用尾巴抓起粉笔，并在黑板上写着点什么。

同学们依旧嬉笑着。铃木和子感到不妥，连忙问社团同学。

\twkkai{你听得清那只猫在说什么吗？}

\twkkai{哈哈哈……它只是在喵喵喵地乱叫而已；你没看大家都为它笑得多开心吗？写的东西也像鬼画符一样，真是太……}

铃木和子突然听不见同学们的笑声了。

\twkkai{这是你的使命……使命……使命……！}

那一声声呵斥在耳边回响着；铃木和子仿佛能从座位上，透过猫的圆形瞳孔，看见里面站着的一个人。惊慌的她站起了身，离开了座位，往后门跑了出去。

不出她所料，猫像猎豹一般追了上来。平时孱弱的她多么希望有人能背着她跑。

\twkkai{是啊……如果他在，该多好啊……}

突然，就在这一刻，她发现了自己心中的情愫；她很快便熟练地用恐惧压制住了它。

……

\twkkai{……总之，这张卡纸上的Landau定理，是一次让假设自己把自己卡死的绝佳范例……}

伊藤佐一屏气凝神地听着，生怕自己泄露了偷窥教授秘密字迹的“秘密”。可是他此刻十分确信，教授的时间也是倒退了的。

突然，咖啡馆最远端的一处玻璃门被猛地推开，而后迅速地被锁上。

进来的人正是铃木和子。她气喘吁吁地跑到教授身后，惊讶地和伊藤佐一四目相对。

\twkkai{铃木同学，你怎么也在这里……难道你的时间也错乱了吗？}

铃木同学气喘得根本接不上话，只好指着门外——那只猫正用爪子疯狂地刮蹭着门玻璃，发出难听的“吱吱”声。

\twkkai{教授……今天可是星期一呀……你怎么没来给我们上课……}铃木和子上气不接下气地说着。

\twkkai{是吗？}教授颇显疑惑。

\twkkai{不好；}伊藤佐一开始焦虑起来。他想拉着铃木和子逃跑，但他不敢。

\twkkai{哦，对了——这是你的吗？}

铃木和子把奖章递给了伊藤佐一。

教授全程都没有留意到他们的舞台剧，只是默默地看回那张卡纸。

伊藤佐一接过奖章，突然意识到什么。他突然变得淡定下来，清了清嗓子，便示意让铃木和子坐下。铃木和子瞪了瞪伊藤佐一。

伊藤佐一依旧招着手，让铃木和子坐在教授身旁。

\twkkai{教授，可以帮我们看看这个吗——就当是给我们辅导了；}伊藤佐一递出奖章。

猫抓得更起劲了。但抓着抓着，它开始疲惫，便潜伏在门边。

而无精打采的教授一开始极不情愿地看着奖章，不以为意。但很快，他的眼神愈发显露出光彩。

\twkkai{你问得也太是时候了；伊藤同学，你不觉得，刚才发生的一切是那么地蹊跷吗？这就像这个命题的证明——你需要先把所有凌乱的拼图拼在一起，才能找到丢失的那一块……}

教授又指了指背后的标语。

\begin{center}
\twkkai{暗处不是光的缺席，是光的另一种形式。}
\end{center}

\twkkai{缺口不是拼图的缺席，而是拼图的另一种形式啊；}教授语重心长地说着，仿佛在暗示伊藤佐一什么。

\twkkai{所以呢，第一步，就是拼好拼图；对于我们来说，这一步具体就是……}

\twkkai{
定义  
\[
\zeta_m(s)=\prod_{\chi\in\widehat{G}}L(s,\chi),
\]
其中 \(G=(\mathbb{Z}/m\mathbb{Z})^\times\)，乘积遍历模 \(m\) 的所有特征。
}\index{zeta@$\zeta_m(s)$（特征 $L$-函数之积）}
\index{zeta@$\zeta_m(s)$（特征 $L$-函数之积）!欧拉乘积表达式}

教授的手优雅地在卡纸背面画着婉转的曲线，像快艇划破水面般写下精确的算式。

\twkkai{把所有包含这些指标的\(L\)-函数乘起来，得到一个新的、有序的整体——我们的拼图便拼好了。在这里，缺口就等同于有其中一个\(L\)-函数在\(s=1\)的时候等于0——利用的是反证法；}\index{ly@$L(1,\chi)\neq 0$（非平凡特征非零性）!反证法}教授用笔尖指着卡纸上的那个函数。

咖啡馆内的音响突然传来一声声协奏曲的巨响；猫早已被吓跑了。

而铃木和子在猫消失之后，勉强能把注意力摆在教授所画的鬼画符上；她竟然在这一刻第一次真正感受到了数论的魅力。

\twkkai{更妙的还在后头……}

教授打了个喷嚏，颤动着身体。

\twkkai{……伊藤同学，还记得我们第一堂课讲的是什么吗？}

\twkkai{记得；讲的是\(L\)-函数的欧拉乘积分解，条件是\(\Re(s)>1\)；}\index{ou@欧拉（L.\ Euler）}\index{ou@欧拉（L.\ Euler）!欧拉乘积}伊藤佐一也早已彻底地把猫抛于脑后。而现在，趁铃木和子在场，他不想浪费每一次机会去表现自己。

但铃木和子对他说的话表现得一头雾水。

\twkkai{很好；我们拼好这块拼图的目的，就是为了对这个全新的\(\zeta_m(s)\)进行分解……那我们就可以写下这些……}

教授不紧不慢地写着；而伊藤佐一仍浸润在兴奋之中，全然忘记了自己所即将再次身处的危机。

\twkkai{
对 \(\Re(s)>1\)，利用 \(L\)-级数的欧拉乘积得到\[\zeta_m(s)=\prod_{\substack{p\in P\\ p\nmid m}}\left(\frac{1}{1-p^{-s\cdot \operatorname{ord}(p)}}\right)^{\frac{\phi(m)}{\operatorname{ord}(p)}},\]  
其中 \(\operatorname{ord}(p)\) 是 \(p\) 在 \(G\) 中的阶。

展开每个因子，可知 \(\zeta_m(s)\) 可写为一个系数非负的狄利克雷级数  
\[
\zeta_m(s)=\sum_{n=1}^{\infty}\frac{a_n}{n^s},\qquad a_n\ge 0.
\]
}\index{zeta@$\zeta_m(s)$（特征 $L$-函数之积）!系数非负 $a_n\geq 0$}
\ \ \quad \twkkai{等等——为什么会有这一切的发生呢？我是指，这些欧拉乘积分解的推导过程……；}铃木和子胆怯地问着。

\twkkai{很好，很好——两位同学都很积极嘛；我就来给你们解释清楚。这个分解的核心思路是：交换乘积次序，利用特征在元素上取值的均匀性，化特征乘积为单位根\index{dan@单位根}的乘积。其实就是这样的：
\[
\zeta_m(s)=\prod_{\chi\in\widehat{G}}L(s,\chi)=\prod_{\chi\in\widehat{G}}\prod_{p\,\nmid\, m}\left(1-\frac{\chi(p)}{p^{s}}\right)^{-1}=\prod_{p\,\nmid\, m}\prod_{\chi\in\widehat{G}}\left(1-\frac{\chi(p)}{p^{s}}\right)^{-1}.
\]

现在我们不妨来考虑一下内层乘积的倒数
\[
P_p(s):=\prod_{\chi\in\widehat{G}}\left(1-\frac{\chi(p)}{p^{s}}\right).
\]

那么，针对这个特征群\(\widehat{G}\)里的特征\(\chi\)——模 \(m\) 下素数 \(p\) 的周期\(\operatorname{ord}(p)\)早就已经把它们的取值锁死在 \(\operatorname{ord}(p)\) 次单位根里，而全体特征刚好把这些根一个不落地均匀瓜分，恰好都现身 \(\phi(m)/f\) 次。\footnote{设 \(G=(\mathbb{Z}/m\mathbb{Z})^{\times}\)，\(|G|=\phi(m)\)，素数 \(p\nmid m\) 在 \(G\) 中的阶为 \(f=\operatorname{ord}(p)\)。特征群 \(\widehat{G}\) 限制到循环子群 \(\langle p\rangle\) 上为满射，且每个纤维大小都等于 \(\phi(m)/f\)，所以当 \(\chi\) 跑遍 \(\widehat{G}\) 时，\(\chi(p)\) 恰好取遍所有 \(f\) 次单位根，每个根重复 \(\phi(m)/f\) 次。}

于是
\[
\prod_{\chi}\left(1-\frac{\chi(p)}{p^{s}}\right)
= \Bigl[\prod_{\zeta^{f}=1}\bigl(1-\zeta\,p^{-s}\bigr)\Bigr]^{\phi(m)/f}.
\]
利用恒等式 \(\prod_{\zeta^{f}=1}(1-\zeta X)=1-X^{f}\)，代入 \(X=p^{-s}\) 得
\[
\prod_{\chi}\left(1-\frac{\chi(p)}{p^{s}}\right) = \bigl(1-p^{-fs}\bigr)^{\phi(m)/f}\dots
\]}\index{dan@单位根!恒等式 $\prod_{\zeta^f=1}(1-\zeta X)=1-X^f$}
\ \ \quad \twkkai{教授——我想打断一下；这里为什么要对内层乘积取倒数呢？}伊藤佐一这次彻底忘记了自己想吸引铃木和子的意图，开始为眼前的推导感到真正的烦恼和不解。

\twkkai{就是为了方便运算……}

\twkkai{这个理由似乎不太充分呀……；}伊藤佐一忐忑地说着。他不敢知道但是又急切地想知道教授怎么回应。

\twkkai{伊藤同学，你说得对，“为了方便”不是理由，只是借口；真正的理由是——我们要让某种对称性自己浮出水面。}

伊藤佐一和铃木和子都停了下来。

\twkkai{如果按照原来的内层乘积，保留倒数的话……那么，\(\chi(p)\)对\(p^{-s}\)的作用就被彻底埋没了，它的精准击中\(\operatorname{ord}(p)\)次单位根的性质便不容易被我们发现了；}教授拿起笔，在那几个\(\chi(p)\)和\(p^{-s}\)符号上反复圈点着。

\twkkai{虽然，不取倒数的话，也可以算，也可以利用恒等式的另一种形式；但是，这样便看不见特征\(\chi(p)\)排列的均匀程度了——我可不是危言耸听；}教授推了推眼镜，笑了笑。

\twkkai{这就是取倒数的真正动机——结构性的，而不是技术性的。我们做数学，永远是为了前者——为了看清结构。}

教授停了下来，看着伊藤佐一。\twkkai{你刚才的问题很好，问到我了。我的确没有仔细考虑到这一点。你以后做数学，永远不要只满足于“方便”二字；}他又看了看铃木和子。

伊藤佐一则停了下来，感受着。他听懂了，更听出了教授口中的深思熟虑和认真的态度。他能看出，教授年轻时，也似乎被这些问题卡住过很久。

教授又拿起了刚放下的笔。

\twkkai{
好了；那我们回到正题吧。把上面的结果代回对 \(p\) 的连乘积：
\[\boxed{
\zeta_m(s)=\prod_{p\,\nmid\, m}\left(\frac{1}{1-p^{-\operatorname{ord}(p)s}}\right)^{\frac{\phi(m)}{\operatorname{ord}(p)}}}.
\]
}

教授圈点着什么；周围依旧喧闹。

\twkkai{
因为 \(f=\operatorname{ord}(p)\) 整除 \(\phi(m)\)，指数 \(r:=\phi(m)/f\) 是正整数。利用二项式展开
\[
(1-X)^{-r}=\sum_{k=0}^{\infty}\binom{r+k-1}{k}X^{k},
\]
每个局部因子可写成
\[
\sum_{k=0}^{\infty}\binom{\frac{\phi(m)}{f}+k-1}{k}p^{-kfs}.
\]
}

教授又开始飞快地讲授了，不留伊藤佐一和铃木和子思考的时间和空间。

\twkkai{
其系数非负，且仅当 \(n\) 是 \(p^{f}\) 的幂时才会出现非零项。

所有素数局部因子的乘积（在 \(\Re(s)>1\) 时绝对收敛）展开后，就得到一个狄利克雷级数
\[
\zeta_m(s)=\sum_{n=1}^{\infty}\frac{a_n}{n^{s}},
\]
其中每个 \(a_n\) 是有限个非负二项式系数的乘积求和，自然满足 \(a_n\ge0\)。
}\index{dilikeleijishu@狄利克雷级数!系数非负条件}

\twkkai{好了，现在你们弄清楚了吗？}

这时，又到铃木和子摇头了。

\twkkai{教授，\(a_n\)的显式构造，你可以写出来吗？我没有看到下标\(n\)和模\(m\)之间的联系……}

\twkkai{当然可以；}教授点了点头，似乎对铃木和子带着极高的赞赏。

伊藤佐一看着铃木和子，对她的话语感到敬佩之时，却又同时产生了一丝嫉妒——仿佛教授只应该和他对话；或者说，只有他才配问教授他心中的问题。

一阵寒颤突然惊扰了他——他此刻是那么对不起铃木和子，只因他有这么一个不好的动机：他喜欢她，却同时不希望她有所领悟——仿佛这样会威胁他的追求者地位；简言之，让他变得被动。

教授打断了他的胡思乱想。

\twkkai{这个\(a_n\)的显式形式，就是这样的：
\begin{equation*}
    a_n =
    \begin{cases}
        \displaystyle \prod_{p \mid m} \binom{r_p + k_p - 1}{k_p}, & \text{\twkkai{如果} } n \text{ \twkkai{可唯一写成 }} n = \prod_{p \mid m} p^{k_p},\ k_p \ge 0, \\[6pt]
        0, & \text{\twkkai{否则}.}
    \end{cases}
\end{equation*}
}

教授的字迹是那么优雅——他用笔在卡纸上划出的大括号和两条圆括号，那弧线是如此美妙婉转——以至于一旁的两个人都屏息凝神地看着，不敢眨眼。

\twkkai{好了，现在就是对付那该死的猫的时候了……；}教授左右张望着。伊藤佐一和铃木和子也跟着左右张望着。

\twkkai{也是对付这个\(L(1,\chi)\)的非零性的时候了……这叫以其人之道还治其人之身；}教授邪魅一笑，又开始飞速写了起来。

\twkkai{
如果我们假设——存在某个非平凡特征 \(\chi\neq 1\)，使得 \(L(1,\chi)=0\)；另外\(L(s,1)\) 在 \(s=1\) 处有单极点（来自 \(\zeta(s)\) 的极点），而非平凡的 \(L(s,\chi)\) 在 \(\Re(s)>0\) 内全纯。  

那么，\(\zeta_m(s)=\prod_{\chi}L(s,\chi)\) 在 \(s=1\) 处极点便会与零点抵消，从而在 \(s=1\) 全纯。这是因为，在复变函数中，极点和零点之所以能“抵消”，是因为它们在函数相乘时，对奇点的影响“恰好相反”。具体来说，是这样的：

如果\(f(s)\) 在 \(s=s_0\) 有 \(m\) 阶极点，那么在其附近 \(f(s) \approx \dfrac{A}{(s-s_0)^m}\)（\(A\neq 0\)）。
而若 \(g(s)\) 在 \(s=s_0\) 有 \(n\) 阶零点，那么 \(g(s) \approx B\,(s-s_0)^n\)（\(B\neq 0\)）。}\index{ji@极点}\index{ji@极点!$\zeta_m(s)$ 在 $s=1$ 的正则性（反证法中的全纯性）}\index{ji@极点!单极点（$L(s,1)$ 在 $s=1$ 处）}\index{ji@极点!零点阶数与极点阶数的抵消}

教授推了推眼镜，又看了看伊藤佐一和铃木和子，咳嗽了几声，又继续说着。

\twkkai{好了，现在，把它们乘起来就会有：
\[
f(s)g(s) \approx \frac{A B}{(s-s_0)^m} \cdot (s-s_0)^n = A B \cdot (s-s_0)^{n-m}.
\]

如果\(n = m\)（零点阶数等于极点阶数），那么乘积在 \(s_0\) 处就是解析且非零的了。}

\twkkai{这就好比极点和零点这对冤家——你们应该是谈过恋爱的吧——如果它们也碰在一起谈恋爱，就会互相抵消掉彼此的极端。我接下来要讲的证明就是要告诉你们：在狄利克雷的世界里，这种救赎关系啊——是不存在的。}教授漫不经心地说着，用话语点缀着咖啡馆。

伊藤佐一的内心泛起了不安；而他和铃木和子的脸，同时“唰”地红了起来。

\twkkai{那么，由Theorem \ref{landau} (Landau定理)……}

教授把卡纸翻过来，用手指敲了敲。

\twkkai{伊藤同学，这条定理你应该很熟了吧？}

伊藤佐一点了点头。

\twkkai{好；那以后铃木同学若是不懂，就找他了喔；}教授又扭头看向铃木和子；铃木和子羞涩地低下头，玩弄着手指。

教授又把卡纸翻了过来。

\twkkai{很好；有这条定理之后，就可以知道——\(\zeta_m(s)\) 的级数表达式将在 \(\Re(s)>0\) 上收敛。我们还差最后一步！}

教授仿佛在给自己开庆功宴，猛地啜饮了几口拿铁。

\twkkai{
我们现在需要来考察 \(\zeta_m(s)\) 乘积表达式中一个因子的下界： 
\[
(1-p^{-\operatorname{ord}(p)s})^{-\frac{\phi(m)}{\operatorname{ord}(p)}} \ge \frac{1}{1-p^{-\phi(m)s}}=\sum_{i=0}^{\infty} p^{-i\phi(m)s},
\]}\index{zeta@$\zeta_m(s)$（特征 $L$-函数之积）!下界估计}
\twkkai{这个不等式其实就等价于
\(\frac{1}{(1-x)^a}\geq\frac{1}{1-x^a},\text{\twkkai{对于}}0<x<1,\ a\in\mathbb{N}.\)
}

\twkkai{
因此  
\[
\zeta_m(s) \ge \prod_{p\nmid m}\left(\sum_{i=0}^{\infty}p^{-i\phi(m)s}\right)=\sum_{(n,m)=1}\frac{1}{n^{\phi(m)s}}.
\]  
}

伊藤佐一瞪大了他惊奇的眼睛，聚精会神地继续听着。

\twkkai{
令 \(s=\frac{1}{\phi(m)}\)，右端级数 \(\sum_{(n,m)=1} \frac{1}{n}\) 发散，这与 \(\zeta_m(s)\) 在该处收敛矛盾。  

结论：假设不成立，所有非平凡特征的 \(L(1,\chi)\neq 0\)。}\index{ly@$L(1,\chi)\neq 0$（非平凡特征非零性）!令 $s=1/\phi(m)$ 导出矛盾}

一阵沉寂。

\twkkai{伊藤同学，铃木同学——你们没听懂吗？}

\twkkai{不不……是教授你讲得太精彩了；}伊藤佐一把自己的手搭在教授的手上，十分恭敬地说着。

\twkkai{好啦，好啦——不用拍我的马屁……；}教授放下了他的钢笔。

\twkkai{教授，我没有这个意思……我是第一次感到如此地震惊；}伊藤佐一挺直了身板，仿佛随时准备上台演讲一般，在铃木和子面前假装摆好了架势。

\twkkai{那很好；不过，你的震惊，就先留到之后在办公室再和我表达吧——今天也是累坏了……回去还要收拾烂摊子呢；}教授拿起拿铁，站了起来。

\twkkai{好……教授慢走；}伊藤佐一说着。

\twkkai{……慢走……；}铃木和子也配合地说着；她的脸依旧泛着樱花红。

教授离开了；咖啡馆的玻璃门依旧原地摆荡着。伊藤佐一和铃木和子，竟然同时坐了下来。

\twkkai{谢谢你每次都费心还我奖章，真的不好意思。需要我……请你一杯吗？}伊藤佐一的手在大腿上依次起舞。

这时，咖啡馆外的喷泉，也依次地向上泵出水柱，像竖琴的银弦一般，随音乐交错变换着；地面传来阵阵哗啦声。

铃木和子摆了摆手。

\twkkai{不用了。不过，我想确认一下——今天是星期一吗？}

\twkkai{今天不是星期天吗？}

此刻，铃木和子再次犹疑了起来。她不确定他这次有没有感受到周遭世界的怪异——因为他自己仿佛就身在怪异之中。

而这时，那只猫又悄悄地潜伏在门的一侧，虎视眈眈地盯着两人。铃木和子连忙站起了身，顾不得那么多，便直接拉起了伊藤佐一的手臂。一股暖流涌进伊藤佐一的双手，逐渐蔓延到他的双耳。

\twkkai{今天根本不是星期天……快走吧——有不好的事要发生了；}铃木和子回头瞟了伊藤佐一一眼，又极为小声地补充了一句。

\twkkai{……记得藏好你的奖章……}

伊藤佐一的手心出汗了；可是铃木和子依旧紧紧地只是拽着他的手臂，仿佛不愿和他的手相遇。

\twkkai{不过也没关系了——还是快点跟着她跑吧……}


